\section{Bhatt--Lurie morphisms retaining the full support lattice}
\label{sec:BLM}

\paragraph{Notation, sources, and scope.}
The user's notation is retained: $Z_0$ denotes absence, $Z_1$
ungraded presence represented by $\tau$, and $Z_2$ the even/odd
parity carried by integer elements. No $Z_2$ parity is assigned
to $\tau$. The coefficient ring of $p$-adic integers is written
$\mathbf Z_p$; it is not the user's state label $Z_2$.
The original split-support absorber is written $\tau_{\rm old}$,
because its multiplication differs from that of the new unit.

\begingroup\raggedright
The primary sources are Bhargav Bhatt and Jacob Lurie,
\href{https://arxiv.org/abs/2201.06120v1}{\emph{Absolute Prismatic
Cohomology}}, original \nolinkurl{apc.tex}, Definition
\texttt{definition:\allowbreak{}generalized-\allowbreak{}Cartier-\allowbreak{}divisor} (2057--2065),
Definition \texttt{definition:\allowbreak{}Cartier-\allowbreak{}Witt-\allowbreak{}divisor}
(2086--2095).\par
Remark \texttt{remark:\allowbreak{}arithmetic-\allowbreak{}lambda-\allowbreak{}line} (2108--2120),
Construction \texttt{construction:\allowbreak{}Cartier-\allowbreak{}Witt-\allowbreak{}atlas} (2149--2157),
and Remark \texttt{remark:\allowbreak{}WCart-\allowbreak{}lift-\allowbreak{}Frobenius} (2952--2959);
and their
\href{https://arxiv.org/abs/2201.06124v1}{\emph{The prismatization
of $p$-adic formal schemes}}, Construction
\texttt{GenCartDivAnim} (286--301), Definition
\texttt{DefAnimPrism} (318--343), and Example
\texttt{CartWitttoAnim} (386--389).
The local originals were matched byte for byte against the retained
arXiv source archives. Reading these stated passages is not a claim
to have read the complete 242-page first work.\par\endgroup

The support carrier is retained from \emph{Split Support Geometry},
version 11, source byline The Clankers,
\href{https://github.com/KokunoYumeto/zeta-function-research-reader/blob/a2851f9eb352e7e4376bc629de86fa4ed9c1c434/workbenches/splitzero-tandem/foundations/split-support-geometry-v11/split_support_geometry_arithmetic_curve_v11.tex#L1226}
{Definition \texttt{def:lattice-split}}. The prime-character
representation is reproved in FSI1--3 in the accompanying source.
The construction below is a morphism into the stated Cartier--Witt
moduli problem, not an assertion that an arbitrary semiring is
already a prism or that the new unit is the historical absorber.

\subsection{The affine-line morphism with its complete input}
For a commutative ring $R$, the groupoid
$[\mathbf A^1/\mathbf G_m](R)$ classifies pairs $(J,\beta)$ with
$J$ an invertible $R$-module and $\beta:J\to R$ linear.
A morphism to $(J',\beta')$ is an isomorphism $u:J\to J'$
such that $\beta=\beta'u$. For free lines with coordinates
$a,b\in R$, this condition is
\begin{equation}\tag{BLM1}
u\in R^\times,\qquad a=bu.
\end{equation}
The inverse morphism has coordinate $u^{-1}$; composition is
multiplication, which preserves BLM1. In particular the
automorphism group of $(R,a)$ is
\begin{equation}\tag{BLM2}
\{u\in R^\times:a(u-1)=0\}.
\end{equation}
At $a=1$ it is trivial, and at $a=0$ it is $R^\times$.
The zero map from a line is therefore not the absent line.
This is an exact property of the moduli object, with no inference
that the user's $Z_0$ is a line or that $\tau$ carries parity.

Fix a prime $p$ and a ring $R$ on which $p$ is nilpotent.
Write $W(R)$ for its $p$-typical Witt ring and $w_0:W(R)\to R$
for its first-coordinate homomorphism. A Cartier--Witt divisor
$(I,\alpha)$ has $I$ invertible over $W(R)$, nilpotent image of
$w_0\alpha$, and unit ideal generated by $\delta\alpha(I)$.
Bhatt--Lurie's affine-line morphism is explicitly
\begin{equation}\tag{BLM3}
\begin{split}
\mu_R:\mathrm{WCart}(R)&\longrightarrow
 [\widehat{\mathbf A}^1/\mathbf G_m](R),\\
(I,\alpha)&\longmapsto
\left(R\otimes_{W(R)}I,\ r\otimes x\mapsto r\,w_0(\alpha x)\right).
\end{split}
\end{equation}
The formula is balanced because $w_0$ is a ring homomorphism.
Tensoring an isomorphism $u$ of input lines gives an isomorphism
of output lines and preserves $\alpha=\alpha'u$, so this is a
functor on the stated groupoids. For a ring map $R\to S$, extension
of scalars along $W(R)\to W(S)$ followed by $w_0$ agrees with
extension along $R\to S$. The elementary tensor associativity
isomorphism identifies both outputs and their maps. This proves
the naturality square. The image section is nilpotent by the first
Cartier--Witt condition; hence the formal completion is retained.
In the arXiv source, the direction printed in Remark
\texttt{remark:\allowbreak{}functoriality-\allowbreak{}WCart}, lines 2127--2129, is reversed
relative to extension of scalars; the tensor formula just given
fixes the direction explicitly, consistently with lines 2138--2140.

\subsection{All support labels enter a compatible prismatic family}
Let $L$ be a nontrivial bounded distributive lattice. Let $X$ be
the set of all proper prime lattice ideals and define
$\chi_\lambda(P)=0$ for $\lambda\in P$ and $1$ otherwise.
Prime separation implies that $\lambda\mapsto\chi_\lambda$ is
injective. Here is its separation argument: if $a\not\le b$,
maximize an ideal containing $b$ and disjoint from the filter
$\uparrow a$ by Zorn's lemma. If $x,y$ both lie outside that
ideal but $x\wedge y$ lies in it, maximality gives
$a\le u\vee x$ and $a\le v\vee y$ with $u,v$ in the ideal.
Distributivity then implies $a\le u\vee v\vee(x\wedge y)$,
a contradiction. Thus that ideal is prime and separates $a,b$.
The full characteristic identities, valid also in characteristic
$p$ with the displayed subtraction retained, are
\begin{equation}\tag{BLM4}
\chi_{0_L}=0,\quad \chi_{1_L}=1,\quad
\chi_{\lambda\wedge\nu}=\chi_\lambda\chi_\nu,\quad
\chi_{\lambda\vee\nu}=\chi_\lambda+\chi_\nu-\chi_\lambda\chi_\nu.
\end{equation}
They follow at each prime from its join and meet conditions.

For every integer $N\ge2$, put
\begin{equation}\tag{BLM5}
B=\operatorname{Map}(X,\mathbf F_p),\qquad
R_N=B[T]/(T^N),\qquad a_{\lambda,N}=T\chi_\lambda.
\end{equation}
All prime coordinates of $L$ occur simultaneously. Polynomial
coefficients in $R_N$ are unique in degrees $0,\ldots,N-1$,
so distinct $\chi_\lambda$ give distinct $a_{\lambda,N}$.
The exact support formulas are
\begin{equation}\tag{BLM6}
\begin{split}
a_{\lambda\vee\nu,N}&=a_{\lambda,N}+a_{\nu,N}
                         -a_{\lambda\wedge\nu,N},\\
a_{\lambda,N}a_{\nu,N}&=T\,a_{\lambda\wedge\nu,N}.
\end{split}
\end{equation}
These are obtained by multiplying BLM4 by $T$ and $T^2$,
respectively. In particular the extra factor $T$ is not discarded.
This is not a claim that $\lambda\mapsto a_{\lambda,N}$ is
a meet-to-product homomorphism.

Define a free-line Cartier--Witt object by
\begin{equation}\tag{BLM7}
d_{\lambda,N}=[a_{\lambda,N}]+V(1),\qquad
\alpha_{\lambda,N}:W(R_N)\xrightarrow{\,d_{\lambda,N}\,}W(R_N).
\end{equation}
The bracket denotes the Teichm\"uller map and $V$ the Verschiebung.
The Witt coordinates of this element are exactly
$(a_{\lambda,N},1,0,0,\ldots)$. To check this without cancelling
in a torsion ring, check the ghost polynomials over the universal
torsion-free polynomial ring: both expressions have zeroth ghost
coordinate $a$ and $n$th coordinate $a^{p^n}+p$ for $n\ge1$.
Ghost injectivity there identifies their universal coordinates;
specialization gives the equality over $R_N$.
Its first coordinate is nilpotent and its second is a unit.
Thus it lies in Bhatt--Lurie's Cartier--Witt atlas and BLM7
satisfies their divisor definition. Their morphism gives
\begin{equation}\tag{BLM8}
\mu_{R_N}(W(R_N),\alpha_{\lambda,N})
=(R_N,\ m_{T\chi_\lambda}).
\end{equation}
This calculates the image, rather than identifying a label with
an unspecified point of an affine line.

For $2\le M\le N$, the quotient $R_N\to R_M$ sends
$a_{\lambda,N}$ to $a_{\lambda,M}$. Witt functoriality sends
$d_{\lambda,N}$ to $d_{\lambda,M}$, because it preserves
Teichm\"uller representatives and $V$. Tensoring BLM7 therefore
gives BLM7 at level $M$, including the map on its line. We obtain
a compatible formal family at every support label.

\subsection{Exact separation and automorphisms after the quotient}
If the two outputs in BLM8 were isomorphic, BLM1 would give
$T\chi_\lambda=T\chi_\nu u$ with $u\in R_N^\times$.
At each coordinate $P\in X$, the constant coefficient of $u(P)$
is nonzero. Comparing coefficients of $T$ shows
$\chi_\lambda(P)=\chi_\nu(P)u_0(P)$. This is impossible if
exactly one of the two characteristic values is zero. Thus the
two supports coincide and prime separation gives $\lambda=\nu$.
Conversely equal labels give the identity isomorphism. Therefore
\begin{equation}\tag{BLM9}
\lambda\longmapsto[(W(R_N),\alpha_{\lambda,N})]
\quad\hbox{is injective on isomorphism classes.}
\end{equation}
The conclusion for Cartier--Witt objects follows because their
images under the functor $\mu$ are already separated. We do not
call a functor from discrete labels fully faithful: target
automorphisms need not be trivial.

The automorphisms of the output line are calculated by BLM2:
\begin{equation}\tag{BLM10}
\operatorname{Aut}(R_N,m_{T\chi_\lambda})
=\{u\in R_N^\times:T\chi_\lambda(u-1)=0\}.
\end{equation}
On the factor $\chi_\lambda B$, this means
$u=1+T^{N-1}v$ with $v\in\chi_\lambda B$; on the complementary
factor $(1-\chi_\lambda)B$, any unit is permitted. This follows
by comparing all coefficients of $T,\ldots,T^{N-1}$; the
idempotent decomposition of $B$ separates the two factors.
These are the full automorphisms in the affine-line quotient;
no assertion that all of them lift to Cartier--Witt
automorphisms is made.

\subsection{Frobenius moves the retained label into higher order}
On a characteristic-$p$ ring the Witt Frobenius is $W(x\mapsto
x^p)$, and $p=V(1)$. Thus $F$ acts coordinatewise by $p$th
powers, fixes $p$, and preserves the displayed formal family:
\begin{equation}\tag{BLM11}
F^j(d_{\lambda,N})=p+[T^{p^j}\chi_\lambda],\qquad j\ge0.
\end{equation}
Indeed $\chi_\lambda^{p^j}=\chi_\lambda$ at every coordinate,
and $F^j([a])=[a^{p^j}]$. This proves the formula including
the constant $p$ and the complete order $p^j$.
If $p^j\ge N$, every element in BLM11 equals $p$. If $p^j<N$,
their affine-line images $(R_N,m_{T^{p^j}\chi_\lambda})$
are still pairwise nonisomorphic, by the coefficient argument
of BLM9 in degree $p^j$. Hence
\begin{equation}\tag{BLM12}
\begin{cases}
F^j(d_{\lambda,N})=p\text{ for all }\lambda,&p^j\ge N,\\
\lambda\mapsto[F^j(d_{\lambda,N})]\text{ is injective},&p^j<N.
\end{cases}
\end{equation}
For each fixed $j$, the compatible family over all $N$ retains
every label because one can take $N>p^j$. A fixed finite
truncation and the entire formal family therefore have different
Frobenius information, with the loss threshold computed exactly.

\subsection{A nonzero derived class at every support label}
Every BLM7 has the following explicit annihilator element:
\begin{equation}\tag{BLM13}
w_N=[T^{N-1}]\ne0,\qquad d_{\lambda,N}w_N=0.
\end{equation}
Its first Witt coordinate is nonzero, proving $w_N\ne0$.
The identities $[a][b]=[ab]$ and $V(1)[b]=V([b^p])$ give
\begin{equation}\tag{BLM14}
\begin{split}
d_{\lambda,N}w_N
&=[T\chi_\lambda][T^{N-1}]+V(1)[T^{N-1}]\\
&=[T^N\chi_\lambda]+V([T^{p(N-1)}])=0+0=0,
\end{split}
\end{equation}
since $N\ge2$, $p\ge2$, and $p(N-1)\ge N$. All exponents
and summands remain displayed. The Witt product identity used
here follows from $V(x)y=V(xF(y))$; it can be checked on
universal ghost polynomials and specialized as above.

The animated quotient of Bhatt--Lurie is the derived quotient
\begin{equation}\tag{BLM15}
\overline A_{\lambda,N}
=W(R_N)\mathbin{/\mkern-6mu/}d_{\lambda,N},\qquad
K_{\lambda,N}=
[W(R_N)\xrightarrow{d_{\lambda,N}}W(R_N)],
\end{equation}
with the left module in homological degree one. To see its
underlying complex, form the derived tensor product of
$W(R_N)[x]/(x)$ and $W(R_N)$ over $W(R_N)[x]$, where the
second map sends $x$ to $d_{\lambda,N}$. The two-term free
resolution by multiplication by $x$ is exact because $x$ is
a non-zero-divisor in the polynomial ring, regardless of
torsion in its coefficients. Tensoring gives exactly BLM15.
It follows that
\begin{equation}\tag{BLM16}
\pi_0\overline A_{\lambda,N}=W(R_N)/(d_{\lambda,N}),\quad
\pi_1\overline A_{\lambda,N}=\operatorname{Ann}(d_{\lambda,N})
\ni w_N\ne0,\quad \pi_i=0\ (i\ge2).
\end{equation}
The Cartier--Witt source gives an animated prism as in their
Example \texttt{CartWitttoAnim}; the explicit nonzero class
also proves why it cannot be replaced by an injective
classical Cartier divisor. This is a computed derived
contribution, not a proposed future comparison.

The coefficientwise Frobenius defines an actual chain map
\begin{equation}\tag{BLM17}
K_{\lambda,N}\longrightarrow K^{(j)}_{\lambda,N}
=[W(R_N)\xrightarrow{F^j d_{\lambda,N}}W(R_N)],
\qquad (F^j,F^j),
\end{equation}
because $F^j(d x)=F^j(d)F^j(x)$. The witness is sent to
$[T^{p^j(N-1)}]$, which is zero for every $j\ge1$.
Thus this particular nonzero degree-one class is killed by
Frobenius, with both its source and target complexes specified.

\subsection{The return to the original full split-support carrier}
Put $A=W(R_N)$ and retain every pair in
\begin{equation}\tag{BLM18}
G_L(A)=(\{0\}\times L)\cup(A\times\{1_L\}),\quad
(a,\sigma)+(b,\eta)=(a+b,\sigma\vee\eta),\quad
(a,\sigma)(b,\eta)=(ab,\sigma\wedge\eta).
\end{equation}
The scalar maps $\mathbf Z\to A\xrightarrow{w_0}R_N$
induce the exact semiring morphisms
\begin{equation}\tag{BLM19}
G_L(\mathbf Z)\longrightarrow G_L(A)\longrightarrow G_L(R_N),
\qquad (n,\sigma)\longmapsto(n1_A,\sigma)
\longmapsto((n\bmod p)1_{R_N},\sigma).
\end{equation}
They preserve join, meet, amplitudes, global zero and unit by
direct substitution in BLM18. Nonzero amplitudes of an image
still occur only at top support. Every $(0,\sigma)$ is
retained as $(0,\sigma)$, even when a nonzero integer is
annihilated by the last amplitude map. At $p=2$ this is
the arithmetic $Z_2$ receiver with all labels kept.
The first map is injective: evaluation at any $P\in X$ and
$T=0$ gives $A\to W(\mathbf F_p)=\mathbf Z_p$, where the
integer inclusion is injective. Equality of images also
forces equality of support labels.

The new ungraded presence unit acts through the pointed
multiplicative morphism
\begin{equation}\tag{BLM20}
\{0,\tau\}\longrightarrow G_L(A),\qquad
0\longmapsto(0,0_L)=\tau_{\rm old},\quad
\tau\longmapsto(1,1_L).
\end{equation}
All four products of $\{0,\tau\}$ verify this multiplicative map. It is not asserted to preserve the presence addition: the source has $\tau+\tau=\tau$, whereas its receiving images add to $(2,1_L)$. This is an image of $\tau$, not an assignment of parity
to its source. The image acts as identity on every label. This is scalar action on the receiving module; it does not identify $\tau$ with the nilpotent section $T\chi_\lambda$ of the formal affine-line object. Multiplication by zero is a module endomorphism, not an isomorphism of an invertible line.

Lift the differential of BLM15 without dropping supports:
\begin{equation}\tag{BLM21}
D_{\lambda,N}:G_L(A)\to G_L(A),\qquad
D_{\lambda,N}(x,\sigma)=(d_{\lambda,N}x,\sigma).
\end{equation}
It preserves sums and scalar action: both sides of each
identity have the same join or meet label and the same
ordinary linear amplitude. Composition with amplitude
projection is exactly multiplication by $d_{\lambda,N}$.
For the nonzero class in BLM13 it gives
\begin{equation}\tag{BLM22}
D_{\lambda,N}(w_N,1_L)=(0,1_L)=e\ne\tau_{\rm old}.
\end{equation}
The complete amplitude-zero fibre is
\begin{equation}\tag{BLM23}
\{(x,\sigma):d_{\lambda,N}x=0\}=G_L(\operatorname{Ann}d_{\lambda,N}),
\end{equation}
where the right side is the supported additive-module carrier,
with nonzero $x$ allowed only at $1_L$. In contrast,
$D_{\lambda,N}^{-1}(\tau_{\rm old})=\{\tau_{\rm old}\}$
for nontrivial $L$: an output at bottom must have input
label bottom, and the original carrier then forces amplitude
zero. Thus the ordinary global-zero fibre would miss the
computed class, whereas the supported-zero fibre retains it.
Coefficientwise Frobenius also lifts as
$(x,\sigma)\mapsto(F^j x,\sigma)$ and commutes with BLM21
into the lifted BLM17. Even when $F(w_N)=0$, its image is
$(0,1_L)$ and retains its support.

For the one-element lattice the split carrier is simply $A$;
there are no distinct support labels or prime characters to
embed. Taking $R_N=\mathbf F_p[T]/(T^N)$ directly gives the
same arithmetic and annihilator calculation. In that case
$e=\tau_{\rm old}$, so the distinction in BLM22--23 is
explicitly absent, exactly as the original carrier requires.

\subsection{The complete annihilator and its automorphisms}
The witness BLM13 is part of an explicitly computable whole module.
For $q\ge1$, an idempotent $c\in B$, and $N\ge2$, write
\begin{equation}\tag{BLM24}
d_{N,q,c}=p+[T^q c],\qquad h_N=\lceil N/p\rceil.
\end{equation}
For a Witt vector $x=(x_0,x_1,\ldots)$ the exact answer is
\begin{equation}\tag{BLM25}
\begin{split}
d_{N,q,c}x=0
&\ \Longleftrightarrow\ 
x_n\in T^{h_N}R_N\ \hbox{and}\ cT^{qp^n}x_n=0
\quad\hbox{for every }n\ge0,\\
\operatorname{Ann}(d_{N,q,c})
&=\operatorname{Ann}(p)\cap\operatorname{Ann}([T^qc]).
\end{split}
\end{equation}
In particular, on $cR_N$ the $n$th coordinate belongs to
$T^{\max(h_N,N-qp^n,0)}cR_N$; on $(1-c)R_N$ it belongs to
$T^{h_N}(1-c)R_N$. The letter $c$ here denotes an idempotent
coefficient and does not change the supported-zero symbol $e$.

Here is a proof including the possible cancellation between the two
summands of $d$. Frobenius on $R_N$ sends
$\sum b_iT^i$ to $\sum b_iT^{pi}$, because $b_i^p=b_i$ in $B$.
Its kernel is therefore precisely $T^{h_N}R_N$. The identities
\begin{equation}\tag{BLM26}
px=V(Fx),\qquad
[T^qc]x=(T^qcx_0,T^{qp}cx_1,T^{qp^2}cx_2,\ldots)
\end{equation}
and injectivity of $V$ prove both annihilator descriptions on
the right of BLM25. They immediately prove that those conditions
are sufficient for $dx=0$.

For necessity, split the ring by the idempotent $c$. On the
complement $d=p$, already handled. On the $c$ factor write
$d_q=p+[T^q]$. If $q\ge N$, again $d_q=p$. Otherwise the zeroth
coordinate of $d_qx=0$ gives $T^qx_0=0$, so
$x_0\in T^{N-q}R_N$. If $pq<N$, then
$p(N-q)>(p-1)N\ge N$ and hence $x_0^p=0$.
Writing $y=(x_1,x_2,\ldots)$, the exact identities
\begin{equation}\tag{BLM27}
\begin{split}
x&=[x_0]+V(y),\quad [T^q]V(y)=V([T^{pq}]y),\quad
p[x_0]=V([x_0^p]),\\
d_qx&=[T^qx_0]+V([x_0^p])+V(d_{pq}y)=V(d_{pq}y)
\end{split}
\end{equation}
reduce the assertion to the next coordinate. Repetition reaches
$q'<N\le pq'$. At this stage $T^{q'}x_0'=0$ and
$T^{pq'}=0$, so the full vector $[T^{q'}]x'$ vanishes, by
BLM26. The equation now forces $px'=0$, and all remaining
coordinates have $p$th power zero. The preceding steps also
gave the required powers and the coordinate equations
$T^{qp^n}x_n=0$ at each earlier index. This proves necessity
and BLM25 without assuming the two summands cannot cancel.

For a chosen free-line divisor, an automorphism after forgetting
the chosen basis is a Witt unit $v$ satisfying $dv=d$. Thus
\begin{equation}\tag{BLM28}
\operatorname{Aut}(W(R_N),m_{d_{N,q,c}})
=1+\operatorname{Ann}(d_{N,q,c}).
\end{equation}
Every displayed $1+x$ is a unit: its zeroth coordinate is
$1+x_0$ with $x_0\in T^{h_N}R_N$ nilpotent. A Witt vector is
invertible exactly when its zeroth coordinate is invertible;
successive multiplication equations determine its inverse
coordinate by coordinate. For the family BLM7 the element
$1+[T^{N-1}]$ is a nonidentity automorphism. Thus every finite
stage has genuinely nontrivial stabilizer, calculated by BLM25.

At $p=2,N=2$ the whole calculation reads
\begin{equation}\tag{BLM29}
\begin{gathered}
d_{\lambda,2}=2+[T\chi_\lambda],\quad
\operatorname{Ann}(d_{\lambda,2})
=\{(x_n)_{n\ge0}:x_n\in TR_2\ \forall n\},\\
[T]\ne0,\quad d_{\lambda,2}[T]
=[T^2\chi_\lambda]+V([T^2])=0,\quad
F(d_{\lambda,2})=2,\quad F([T])=0.
\end{gathered}
\end{equation}
The coefficient class in degree one is nonzero before Frobenius.
Its lifted differential still ends at the supported zero in BLM22.

\subsection{A fully faithful map into the compatible formal towers}
Fix $j\ge0$. Form the groupoid whose objects are compatible systems
of Cartier--Witt objects over the rings $R_N$, $N\ge2$, with
specified transition isomorphisms along the truncations. A morphism
is a family of isomorphisms commuting with these transitions.
The systems BLM11 have their canonical transitions from BLM7.
There is a functor from the discrete groupoid on the set $L$:
\begin{equation}\tag{BLM30}
\mathcal D_j:L_{\rm disc}\longrightarrow
\varprojlim_{N\ge2}\mathrm{WCart}(R_N),\qquad
\lambda\longmapsto
\bigl(W(R_N),m_{p+[T^{p^j}\chi_\lambda]}\bigr)_N.
\end{equation}
The limit denotes the just specified groupoid of compatible
systems. This functor is fully faithful.

Indeed distinct labels admit no isomorphism because any stage
$N>p^j$ separates them by BLM12. An automorphism of a single
system is represented by compatible units $1+x_N$, with
$x_N\in\operatorname{Ann}(p+[T^{p^j}\chi_\lambda])$.
Fix a Witt-coordinate index $n$. The compatible coefficients
$(x_N)_n$ give an element of $B[[T]]$. Formula BLM25 forces it
to vanish to order at least $\lceil N/p\rceil$ for every $N$.
Every power-series coefficient is therefore zero. This holds
for each Witt coordinate, giving $x_N=0$ for every $N$.
The identity is the only automorphism. It is the only map from
each labelled object to itself in the discrete source as well.
The two Hom sets consequently agree for every pair of labels,
which proves the assertion.

There is an essential domain statement. For $\lambda\ne0_L$,
prime separation provides a $P$ with $\chi_\lambda(P)=1$.
Evaluation at $P$ proves that $T^{p^j}\chi_\lambda$ is not
nilpotent in $B[[T]]$. Therefore the inverse-limit Witt vector
does not define an ordinary Cartier--Witt point over $B[[T]]$
with the nilpotence condition in BLM3. BLM30 is a morphism
into compatible formal towers, exactly as its codomain states.
At every finite stage that nilpotence condition is satisfied.

\paragraph{Calculated scope.}
The resulting arrows reach Bhatt--Lurie's actual moduli
objects and animated quotients, retain every original lattice
label, separate their deformations, and calculate a nonzero
derived class with its Frobenius image. They do not identify
these classes with the programme's arithmetic spectral
cohomology or assign them a Deligne weight. No positivity
statement is inferred from the existence of these morphisms.
